% !PRINTED none
% Abschnitt-opening (part-title) page: no folio is printed on this page in the scan
% (chapter-opening pages in this book suppress the running head/page number).
% Content: "Sechzehnter Abschnitt." / "Kreistheilung." (both letter-spaced, centred),
% a horizontal divider rule, then "§. 167." and its letter-spaced title, followed
% immediately by the start of ordinary body text (no page break before the body begins).
\begin{center}
\gesperrt{Sechzehnter Abschnitt.}

\gesperrt{Kreistheilung.}

\rule{3cm}{0.4pt}

\S.\,167.

\gesperrt{Die Kreistheilungsperioden und die Periodengleichungen.}
\end{center}

Die wichtigsten unter den Abel'schen Gleichungen sind die, von denen die Bestimmung der Einheitswurzeln abhängt, deren elementare Eigenschaften wir schon im zwölften Abschnitt kennen gelernt haben. Wegen der Beziehung zu der geometrischen Aufgabe, die Kreisperipherie in eine bestimmte Anzahl gleicher Theile zu theilen, heissen alle damit zusammenhängenden Gleichungen auch die \gesperrt{Kreistheilungsgleichungen}.

Der Körper $\Omega$, der die als bekannt angesehenen Grössen enthält, ist hier nur der \gesperrt{Körper der rationalen Zahlen}, den wir den Körper $R$ nennen wollen.

Wir wollen uns zunächst mit den $n^{\text{ten}}$ Einheitswurzeln unter der Voraussetzung beschäftigen, dass $n$ eine ungerade Primzahl sei (die Primzahl 2 bietet zu keinen weiteren Fragen Anlass, da die einzigen zweiten Einheitswurzeln, $\pm 1$, in $R$ enthalten sind).

Ausser der rationalen $n^{\text{ten}}$ Einheitswurzel 1 existiren noch $n - 1$ primitive, die durch die transcendenten Ausdrücke
\[
e^{\frac{2\pi i}{n}}, \quad e^{\frac{4\pi i}{n}}, \quad e^{\frac{6\pi i}{n}} \dots e^{\frac{2(n-1)\pi i}{n}}
\]
dargestellt werden können. Wenn wir eine von ihnen mit $r$ bezeichnen, so ist das ganze System auch durch
\weq{1}{r, \ r^2, \ r^3 \dots r^{n-1}}
%CONT>>
